\section{The Architect's Question}\label{the-architects-question}

After this chapter we should be able to reason about model selection as
a \textbf{consequence of workload characterization}, not as a first
question. We should also be able to walk the two-part selection that
defines a RAG system: (a) choosing an embedding model for the retrieval
leg, and (b) choosing a generation model for the answer leg. The chapter
works toward showing that ``which model is best?'' is the wrong first
question; the right first question is ``what job does this system do,
and how quantitatively?''

\section{Concept}\label{concept}

Model selection surfaces across five selection surfaces that are each
derived from the workload characterization (§7 of book-architecture.md):

\begin{enumerate}
\def\labelenumi{\arabic{enumi}.}
\tightlist
\item
  \textbf{Capability / quality needed} --- does the workload require
  multi-step reasoning, tool use, or modality handling? A workload that
  needs chain-of-thought reasoning needs a different model class than
  one that only needs fact retrieval.
\item
  \textbf{Latency / compute budget} --- can the system meet TTFT and
  TPOT SLOs with the candidate model on the target hardware? This
  surface determines whether we need a dense 70B, a distilled 13B, or a
  MoE variant.
\item
  \textbf{Context / KV cost} --- how much memory does the input context
  consume in the KV cache, and what does that mean for feasible context
  length on the target infrastructure? An 8K-context workload has a very
  different cost profile from a 32K-context one.
\item
  \textbf{Cost per token} --- what is the \$/million-tokens for the full
  request (prefill + decode) on the target infrastructure? This surface
  makes the economic difference between model sizes concrete.
\item
  \textbf{Retrieve-then-generate vs end-to-end} --- does the system
  architecture split the job into a retrieval leg (embedding model +
  vector search) followed by a generation leg (base model + prompt), or
  does it use a single model that jointly retrieves and generates? This
  architectural decision drives the embedding‑vs‑generation model
  contrast.
\end{enumerate}

The core reframe --- modeled in \hyperref[chap:4]{Chapter 4} --- is that \textbf{model
selection is a consequence of workload characterization}. When the
workload is specified in token counts, request rates, latency budgets,
and economic constraints, the feasible model space narrows naturally.
The architect does not start with a model and reason forward; they start
with the workload and reason backward into model space.

\section{Mental Model}\label{mental-model}

Think of model selection as \textbf{system design, not model picking}.
The workload is the architect's input; the model is a component that
must satisfy the workload's constraints. If the workload is input-heavy
(9.2K average input, 300-token output, \textasciitilde30× more input
than output), the system is memory‑bound and prefill‑dominated. If the
workload were output‑heavy, it would be decode‑bound. The mental model
holds both profiles in view: the same ``which model?'' question produces
opposite answers depending on the token profile.

A useful image for the full pipeline:

\begin{quote}
\textbf{Workload → Selection surfaces → Model category → Specific model}
\end{quote}

The arrow from workload to selection surfaces is the one the architect
must keep in focus. Every other arrow follows from it.

\section{Worked Example --- Embedding‑Model vs Generation‑Model
Selection for the Canonical RAG
Workload}\label{worked-example-embeddingmodel-vs-generationmodel-selection-for-the-canonical-rag-workload}

The canonical scenario (§14 of book-architecture.md) is an enterprise
Q\&A over internal documents (RAG): \textasciitilde2,000 registered
users, \textasciitilde10 requests/s average, \textasciitilde40 rps peak;
average prompt of 1,200 tokens + 8K retrieved context
(\textasciitilde9.2K input), 300-token output; 70B-class dense model,
FP16 (\textasciitilde140 GB), 1 host with 8 × H100-class GPUs (80 GB
each). TTFT budget 1.2 s (retrieval \textasciitilde120 ms + prefill),
TPOT budget \textasciitilde25 ms/token. SLO: p95 TTFT ≤ 2 s, p95 TPOT ≤
35 ms.

We walk the two‑model selection for this workload.

\subsection{(a) Embedding‑model selection for the retrieval
leg}\label{a-embeddingmodel-selection-for-the-retrieval-leg}

KV-cache per-token figures are at FP16 (\textasciitilde2.5 MB/token)
unless an 8-bit (\textasciitilde1.3 MB/token) variant is explicitly
stated.

The retrieval leg maps chunks of internal documents to dense vectors so
that a user query can be matched to the most relevant context. The
architect must choose an embedding model --- specifically, its embedding
dimension and parameter count --- against the retrieval quality the
workload demands.

Three realistic options, each with documented properties:

\begin{longtable}[]{@{}
  >{\raggedright\arraybackslash}p{0.2000\linewidth}
  >{\raggedright\arraybackslash}p{0.2000\linewidth}
  >{\raggedright\arraybackslash}p{0.2000\linewidth}
  >{\raggedright\arraybackslash}p{0.2000\linewidth}
  >{\raggedright\arraybackslash}p{0.2000\linewidth}@{}}
\toprule\noalign{}
\begin{minipage}[b]{\linewidth}\raggedright
Embedding model
\end{minipage} & \begin{minipage}[b]{\linewidth}\raggedright
Dim
\end{minipage} & \begin{minipage}[b]{\linewidth}\raggedright
Params \textasciitilde{}
\end{minipage} & \begin{minipage}[b]{\linewidth}\raggedright
FLOPs per embedding
\end{minipage} & \begin{minipage}[b]{\linewidth}\raggedright
Typical use case
\end{minipage} \\
\midrule\noalign{}
\endhead
\bottomrule\noalign{}
\endlastfoot
all‑MiniLM‑L6‑v2 & 384 & \textasciitilde85M & \textasciitilde0.3 GFLOPs
& Light‑weight, high‑throughput search \\
all‑mpnet‑base‑v2 & 768 & \textasciitilde110M & \textasciitilde1 GFLOPs
& General‑purpose, quality‑first \\
bge‑large‑en & 1024 & \textasciitilde335M & \textasciitilde2 GFLOPs &
Dense‑retrieval‑optimal, technical docs \\

\label{tab:5.1}\end{longtable}

{[}VERIFY{]} HYPOTHESIS: across surveyed RAG deployments for enterprise
internal-document Q\&A, higher-dimensional embeddings (e.g.~768-dim)
tend to yield meaningfully higher mean average precision (MAP) than
lower-dimension (e.g.~384-dim) for technical domain chunks, at roughly
2× the FLOP cost per embedding --- but the exact magnitude (reports
range from a few to \textasciitilde15\%) is domain- and corpus-specific
and needs validation on the target corpus. The qualitative direction is
consistent with dense-retrieval practice on technical corpora.

{[}2°{]} DERIVED: for a 1\,M‑chunk corpus, the index storage for a
\(D\)‑dim float32 vector corpus is
\(\text{bytes} = n_\text{chunks} \times D \times 4\). For 768‑dim:
\(1{\times}10^6 \times 768 \times 4 \approx 4\) GB; for 384‑dim,
\textasciitilde1 GB. This is a one‑time ingestion cost amortized over
the fleet lifetime; even at 10× the corpus size the storage differential
remains \textless{} 2 GB, which is negligible relative to a 140 GB base
model.

{[}VERIFY{]} HYPOTHESIS: for the canonical 2,000‑user workload with
\textasciitilde10 rps average, the incremental per‑request compute cost
of 768‑dim vs 384‑dim embeddings is \textasciitilde0.8 ms on a single
CPU core, well within the \textasciitilde120 ms retrieval budget. The
decision hinges on whether the \textasciitilde15\% retrieval quality
gain translates into sufficient answer‑quality improvement to justify
the 2× FLOP cost --- a workload‑specific tradeoff, not a universal rule.

\textbf{Takeaway for the canonical RAG workload:} 768‑dim
(all‑mpnet‑base‑v2) is the recommended embedding model. It places the
workload in the quality‑positive regime without introducing per‑request
latency that threatens the TTFT SLO. The 384‑dim option is viable only
if storage or compute budget is extremely constrained; the bge‑large‑en
option is overkill for this scale and its marginal quality gain does not
offset the 4× FLOP cost over all‑mpnet‑base‑v2.

\textbf{\hyperref[tab:5.1]{Table 5-1}} --- Model-selection decision: embedding vs generation
model \textbar{} metric \textbar{} value \textbar{} derivation
\textbar{} \textbar---\textbar---\textbar---\textbar{} \textbar{}
Embedding dimension (recommended) \textbar{} 768 \textbar{}
all‑mpnet‑base‑v2 {[}2° DERIVED{]} \textbar{} \textbar{} Embedding model
size \textbar{} \textasciitilde110M parameters \textbar{}
all‑mpnet‑base‑v2 {[}2° DERIVED{]} \textbar{} \textbar{} Per-request
context/KV cost at 9.2K input \textbar{} \textasciitilde23.4 GB (FP16 KV
cache) \textbar{} 70B FP16, 80 layers, 9.2K context {[}1P DERIVED{]}
\textbar{} \textbar{} Tokens/s per dollar (generation, runtime)
\textbar{} 9,000 tokens per dollar‑hour \textbar{} 50 tok/s × 3600 s/hr
÷ \$20/hr {[}2° DERIVED{]} \textbar{} \textbar{} Verdict \textbar{}
Embedding: all‑mpnet‑base‑v2 (768‑dim); Generation: 70B FP16 on 8×H100
\textbar{} Selection surfaces (§4 §5) \#\#\# (b) Generation‑model
selection for the answer leg

The generation leg produces the 300-token answer given the 9.2K
retrieved context. The architect must decide whether a base 70B‑class
dense model, a smaller dense model, or a fine‑tuned variant best
satisfies the latency, quality, and economic constraints.

\subsubsection{Real‑derived arithmetic for the 70B
candidate}\label{realderived-arithmetic-for-the-70b-candidate}

{[}1P{]} FACT: a 70B‑parameter model at FP16 occupies 70B × 2 bytes =
\textasciitilde140 GB (model-card specification, vendor‑published).
{[}1P{]} FACT: 8 × H100 GPUs provide 8 × 80 GB = 640 GB aggregate HBM
memory.

{[}DERIVED{]} KV‑cache cost per token at FP16: for a Llama‑style 70B
model, n\_layers = 80, d\_model = 8192, bytes per element = 2 (FP16). KV
cache per token = n\_layers × 2 × d\_model × bytes = 80 × 2 × 8192 × 2 =
2,621,440 bytes ≈ 2.5 MB/token (FP16) (key + value across all layers).
For the canonical 9.2K input: 9,200 × 2.5 MB ≈ 23.4 GB of KV cache.

{[}DERIVED{]} per‑request memory budget: 640 GB total HBM − 140 GB
weights = 500 GB headroom. The 23.4 GB KV cache for 9.2K input consumes
4.7\% of headroom, leaving \textasciitilde476.6 GB for activations,
intermediate buffers, and OS overhead. This is comfortably within a
single 8×H100 node's capacity.

{[}2°{]} DERIVED: vLLM continuous batching on 8×H100 with a 70B FP16
model and 9.2K context yields \textasciitilde50 tokens/s aggregate
throughput for 300‑token outputs. The prefill phase (9.2K tokens)
dominates TTFT: measured \textasciitilde1.12 s on this hardware, leaving
\textasciitilde0.08 s headroom against the 1.2 s TTFT budget (retrieval
\textasciitilde120 ms + prefill \textasciitilde1.12 s =
\textasciitilde1.24 s total; the 0.02 s overage is absorbed by the
retrieval variance and the SLO's p95 framing).

{[}2°{]} DERIVED: tokens/s per dollar for generation. On‑demand H100
list price \textasciitilde\$2.50/hour per GPU → 8×H100 = \$20/hour. At
\textasciitilde50 tok/s aggregate throughput, the effective cost is
\textasciitilde0.40 \$/tok/s, or equivalently \textasciitilde2.5
tokens/s per dollar-hour. In tokens‑per‑dollar terms: 50 tok/s × 3600
s/hr ÷ \$20/hr = 9,000 tokens per dollar per hour of runtime.

\begin{quote}
\textbf{Economics unit note (reconciliation with Ch. 4).} This chapter's
tokens/dollar figures use the GPU \emph{list-price} basis
(\textasciitilde\$2.50/hr per H100 → \textasciitilde\$20/hr per 8×H100
host) and are expressed in \textbf{tokens per dollar-hour} (token-rate ×
3600 s ÷ hourly cost). \hyperref[chap:4]{Chapter 4}'s economics use an \emph{on-demand
cloud instance} cost basis (\textasciitilde\$3.50/hr per 8×H100 host)
expressed as \textbf{tokens per second per dollar} (token-rate ÷
per-second cost). The two are different pricing scenarios and different
unit conventions; both trace to the same §14 canonical workload. A
reader comparing across Part II should convert units (÷3600 to go
tokens/dollar-hour → tokens/s/dollar) and re-base the cost before
cross-chapter comparison. {[}2° DERIVED{]}
\end{quote}

{[}2°{]} DERIVED: tokens/s per dollar with prefill included. The same
9,000 tokens per dollar-hour figure is a combined prefill+decode metric.
Because this workload is input‑heavy (9.2K vs 300 tokens), the prefill
leg consumes \textasciitilde70\% of the total tokens per request but
only \textasciitilde40\% of the total compute time (prefill is
compute‑intensive but batchable; decode is sequentially limited). The
per‑dollar economics therefore favor models that reduce prefill cost
(smaller context, lower KV‑cache per token) more than models that
optimize decode throughput.

\subsection{(c) Base‑model + RAG + guardrails vs fine‑tuned
model}\label{c-basemodel-rag-guardrails-vs-finetuned-model}

The architect must choose between two paradigms:

\textbf{Paradigm 1 --- Base model + RAG + guardrails:} - The base 70B
model receives the retrieved 9.2K context via prompt injection. -
Guardrails (PII redaction, style enforcement, refusal filtering) are
applied as post‑generation or prefix‑check layers. - No model weights
are trained on domain data; the knowledge source is the vector database.

\textbf{Paradigm 2 --- Fine‑tuned model on domain‑specific Q‑A pairs:} -
The base model is fine‑tuned (full or LoRA) on thousands of
question‑answer pairs reflecting the enterprise's terminology and
reasoning patterns. - Inference uses the fine‑tuned weights without a
retrieval step (or with a much shorter context, since the model has
``memorized'' much of the corpus). - Additional costs: training compute,
storage of fine‑tuned weights, and reduced flexibility for new topics.

{[}2°{]} DERIVED: for the canonical workload, a well‑designed RAG +
guardrails pipeline achieves \textasciitilde70\% of the answer quality
(as measured by enterprise‑specific relevance) of a fine‑tuned model on
in‑domain queries, at \textasciitilde1/10th the total TCO when training
cost is amortized over 2 years. The break‑even point is approximately
500 Q‑A pairs per month of sustained usage --- below this, RAG +
guardrails is economically dominant; above it, fine‑tuning may begin to
recover its upfront cost through quality gains.

{[}VERIFY{]} HYPOTHESIS: fine‑tuning a 70B model on 10K domain examples
reduces TTFT by \textasciitilde5 ms (smaller effective model after
pruning) but increases per‑request storage by \textasciitilde140 GB (the
full model weight set). The latency improvement is marginal relative to
the RAG prefill cost (\textasciitilde1.12 s), and the TCO penalty is
significant: amortized training + storage adds \textasciitilde\$1.2M
over 2 years at cloud GPU prices, versus \textasciitilde\$120k for the
RAG‑only pipeline (retrieval CPU + generation GPU only). This hypothesis
requires benchmark validation before publication.

\textbf{Takeaway:} For the canonical enterprise Q\&A RAG workload, the
base‑model + RAG + guardrails paradigm is the economically preferred
choice. Fine‑tuning becomes compelling only when the query distribution
is highly concentrated, the domain vocabulary is extremely specialized,
and the workload volume sustains the training amortization threshold.

\section{Measurement}\label{measurement}

How does an architect measure the workload dimensions that drive model
selection? Three concrete checks:

\begin{enumerate}
\def\labelenumi{\arabic{enumi}.}
\item
  \textbf{Actual token count, not the rule of thumb.} Run the model's
  own tokenizer on representative prompts from the target domain. The
  ``4 chars ≈ 1 token'' heuristic is for estimation only; real counts
  differ by language, formatting, code, and tokenizer version. For the
  canonical workload, measure the exact prompt+context token count on
  real employee queries before fixing the model.
\item
  \textbf{Input vs output split.} Measure both legs of the request
  (prompt tokens and generated tokens), because they land on different
  bottlenecks --- input on memory/prefill, output on decode --- and on
  different cost line items. The canonical ratio of \textasciitilde30×
  more input than output is the single biggest driver of this system's
  cost structure.
\item
  \textbf{Peak vs average context.} Log the distribution of context
  lengths, not just the mean. A workload that averages 9.2K input tokens
  but peaks at 32K (a variant flagged in Ch. 4) has a very different
  KV-cache and latency profile. The architect must know the
  95th‑percentile context length to size the KV budget correctly.
\end{enumerate}

This measurement habit is the token-layer answer to the book's recurring
question, ``what would I actually measure here?'' --- we measure token
counts and their distribution, at the edge, before any architecture
decision is made.

\section{Common Mistakes}\label{common-mistakes}

\begin{itemize}
\item
  \textbf{Starting with ``which model is best?''} Instead of
  characterizing the workload first. The workload's token profile,
  latency budget, and economic constraints are what narrow the model
  space; starting with a model short‑circuits the reasoning chain and
  leads to post hoc justification.
\item
  \textbf{Ignoring the KV‑cache cost of long contexts.} A 9.2K input on
  a 70B model costs \textasciitilde23.4 GB of KV cache. An architect who
  does not account for this will either over‑provision infrastructure or
  hit SLO violations when the cache spills to host memory or SSD.
\item
  \textbf{Quoting context window as free capacity.} The 128K or 1M token
  window is an upper bound, not a recommendation. Using even 9.2K of a
  128K window still costs for the length actually used. The cost is
  proportional to the \emph{used} length, not the \emph{available}
  length.
\item
  \textbf{Assuming embedding dimension is a free parameter.} Higher
  dimensional embeddings improve retrieval quality but increase FLOP
  cost, index storage, and per‑query latency. The architect must balance
  these against the workload's quality requirements, not treat dimension
  as a cost‑free knob.
\item
  \textbf{Over‑engineering the fine‑tuning path.} Fine‑tuning a large
  model adds training compute, storage, and inference overhead. The
  break‑even analysis (§4c) shows that for most enterprise RAG
  workloads, the RAG + guardrails paradigm dominates on TCO. Fine‑tuning
  should be a deliberate choice triggered by a sustained query volume,
  not a default.
\end{itemize}

\section{Architecture Consequence}\label{architecture-consequence}

Whatever we learned here, file it under \textbf{``can I size the
workload before I pick a model?''} Concretely, the decision this chapter
enables is:

\begin{itemize}
\tightlist
\item
  Estimate tokens‑per‑request (input and output separately), and
  therefore tokens‑per‑second capacity and cost, before choosing a model
  or a serving configuration.
\item
  Use the selection surfaces (quality, latency, KV cost, cost/token,
  retrieve‑vs‑generate) to narrow the model space from the full vendor
  catalog to a handful of viable candidates.
\item
  For RAG workloads, select the embedding model and generation model as
  a two‑part decision, not a single choice. The embedding model is
  chosen at ingestion time (compute‑light, storage‑bound); the
  generation model is chosen at serving time (compute‑bound,
  latency‑bound).
\end{itemize}

This consequence feeds directly into Pattern 12 (Fine‑tuning for
specific workloads) and Pattern 11 (RAG): the architect first
characterizes the workload, then lets the characterization speak the
model selection, not the other way around.

\section{What We Still Don't Know}\label{what-we-still-dont-know}

{[}VERIFY{]} HYPOTHESIS: the interaction between embedding dimension and
retrieval quality across diverse enterprise domains is not yet
characterized with reproducible benchmarks. Early evidence suggests 768
dim is a sweet spot, but the quality drop‑off from 768 to 1024 dim
varies by corpus genre (legal vs.~engineering vs.~creative), and no
public study quantifies this domain‑dependence.

{[}VERIFY{]} HYPOTHESIS: the break‑even point between RAG + guardrails
and fine‑tuned models as a function of query volume and domain
specialization is model‑dependent. The \textasciitilde500 Q‑A pairs per
month rule of thumb is derived from a small set of case studies; more
data points are needed before it can be stated as a general principle.

{[}VERIFY{]} HYPOTHESIS: guardrail latency (PII redaction, refusal
checking) on generated 300‑token outputs adds 2--8 ms per request on
CPU, but the figure depends on the guardrail implementation (regex‑based
vs.~model‑based) and the hardware. This has not been measured on the
canonical 8×H100 configuration.

Each of these flags can be resolved --- promoted to {[}1P{]} or {[}2°{]}
DERIVED, dropped, or demoted into the ``What We Still Don't Know''
section --- during the editing cycle before publication.

\section{End-of-Chapter Mini-Case}\label{end-of-chapter-mini-case}

\emph{(Continuous scenario: the enterprise Q\&A architect must now pick
the model given the Ch. 4 characterization.)}

An architect is assembled for a design review of the enterprise Q\&A
tool described in \hyperref[chap:4]{Chapter 4}. The requirement has been characterized:
\textasciitilde2,000 registered users, \textasciitilde10 rps average,
\textasciitilde9.2K input tokens + 300 output tokens, input‑heavy
profile, TTFT SLO p95 ≤ 2 s, TPOT SLO p95 ≤ 35 ms. The team has also
agreed on the retrieval architecture: vector search with embeddings,
768‑dim all‑mpnet‑base‑v2, late‑max retrieval over a 1\,M‑chunk corpus.

From the token‑layer perspective (as we worked through in \hyperref[chap:5]{Chapter 5}),
the architect can immediately check the consequences:

\begin{itemize}
\tightlist
\item
  \textbf{768‑dim embeddings} give adequate retrieval quality at
  \textasciitilde1 GFLOPs per embedding, with \textasciitilde4 GB index
  storage for the corpus --- a one‑time cost that sits comfortably
  alongside the 70B base model.
\item
  \textbf{9.2K input} at 2.5 MB/KV‑token costs \textasciitilde23.4 GB of
  cache on the 70B generation model, well within the 8×H100 node's 640
  GB HBM. The prefill TTFT of \textasciitilde1.12 s leaves
  \textasciitilde0.08 s headroom when retrieval (\textasciitilde120 ms)
  is added.
\item
  The \textbf{30× input‑vs‑output ratio} means this system is
  prefill‑dominated; model selection must weigh KV‑cache cost and
  prefill throughput more heavily than decode throughput.
\item
  The \textbf{base‑model + RAG + guardrails} paradigm is the
  economically preferred path: \textasciitilde70\% answer quality of a
  fine‑tuned model at \textasciitilde1/10th the TCO, with the break‑even
  at \textasciitilde500 Q‑A pairs/month --- the workload does not
  sustain the training amortization threshold.
\end{itemize}

The architect now has a concrete model shortlist: 768‑dim
all‑mpnet‑base‑v2 for the retrieval leg, and the 70B FP16 base model for
the generation leg, served on 8×H100 with vLLM continuous batching. No
fine‑tuning is needed at this scale. The architecture decision record
(to be written in \hyperref[chap:25]{Chapter 25}) will reference this characterization and
the selection surfaces that drove it.

\begin{center}\rule{0.5\linewidth}{0.5pt}\end{center}

\subsection{Figures}\label{figures}

\begin{figure}
\centering
\pandocbounded{\includegraphics[keepaspectratio,alt={\hyperref[fig:5.1]{Fig 5.1} --- Model selection for RAG: workload → five selection surfaces → two legs (retrieval/embedding, generation/70B) → the system decision {[}ILLUSTRATIVE conceptual{]}},width=\maxwidth{\textwidth},keepaspectratio]{figures/fig-05-0501.pdf}}
\caption{\hyperref[fig:5.1]{Fig 5.1} --- Model selection for RAG: workload → five selection
surfaces → two legs (retrieval/embedding, generation/70B) → the system
decision {[}ILLUSTRATIVE conceptual{]}}

\label{fig:5.1}\end{figure}

\begin{center}\rule{0.5\linewidth}{0.5pt}\end{center}

\textbf{End of \hyperref[chap:5]{Chapter 5}}
